\section{工程哲学与方法论}

\subsection{教 Researcher 做 Engineering 比反过来更难}

翁家翌引用同事（一位搞过知名 RL framework 的 PhD）的话：

\begin{importantbox}{Engineering vs Research 的非对称性}
"教一个 researcher 如何做好 engineering，要远比教一个 engineer 如何做好 research 来得难。"在当前 AI 时代，工程能力的价值超过了学术能力。因为 research 直觉可以通过在行业中长期工作积累，但扎实的工程能力需要长时间的刻意训练。
\end{importantbox}

\subsection{迭代速度决定一切}

翁家翌反复强调的核心理念：\textbf{单位时间内的正确迭代次数}是 AI 研究成功的决定性因素。具体来说：

\begin{enumerate}
    \item Infra 的正确性：没有 bug 的 infra 意味着每次实验结果可信
    \item 迭代速度：从一周迭代30次到一周迭代300次，成功率线性增长
    \item Bug 修复吞吐量：单位时间内能修多少 bug，能正确迭代多少次
\end{enumerate}

\begin{warningbox}{算法创新不是瓶颈}
翁家翌认为当前 AI 前沿的瓶颈\textbf{不在算法创新}。"如果你把 bug 全修了，有可能算法连改都不用改就很好。"很多看似需要新算法才能解决的问题，实际上可能只是 infra 的 bug 导致的。这与学术界专注于提出新算法的范式形成了鲜明对比。
\end{warningbox}

\subsection{代码一致性与项目腐化}

翁家翌将\textbf{一致性（consistency）}视为软件项目和组织管理的第一原则。项目腐化的根源是不一致性——多人协作时，每个人的 context 不同，假设无法传递，导致代码膨胀和质量下降。

他认为管公司和管代码库本质上是相同的——都需要保持 consistency。如果不一致，代码库会臃肿，组织架构也会臃肿。解决方案是\textbf{信息流通畅}——上面的决策能无损传达到下面，下面的进展能无损传达到上面。

翁家翌进一步将这个问题推向了极致：\textbf{context sharing 理论上应该由一个拥有无限长 context 的 agent 来替代}。因为人脑的 context 是有限的，无法同时存储整个组织的所有信息。但 AI 可以。未来也许每个公司都有一个这样的无限 context agent 来当 CEO，负责所有的 sharing 和 decision——"可能没有比这样的 agent 更适合的 decision maker 了。"这个展望将他的工程哲学、AI 信念和组织管理思考统一了起来。

\subsection{本章小结}

翁家翌的工程哲学可以概括为三个核心原则：（1）\textbf{一致性优先}——一个人把所有东西全包了，虽然不利于长期扩展，但能保证一致性；（2）\textbf{迭代速度优先}——不追求一次性的天才算法，而是通过快速试错逼近最优解；（3）\textbf{正确性优先}——修 bug 比写新 feature 更重要。


